\documentclass[a4paper,11pt]{article}
\usepackage[utf8]{inputenc}
\usepackage[T1]{fontenc}
\usepackage[brazilian]{babel}
\usepackage{hyperref}

\title{Ferramentas de Linha de Comando: Git, Make e Sed}
\author{Nome do Aluno}
\date{\today}

\begin{document}

\maketitle

\section{Introdução}

Este relatório resume três ferramentas de linha de comando vistas nas semanas iniciais
da disciplina: o controlador de versão \texttt{git}, o automatizador de compilação
\texttt{make} e o editor de fluxo \texttt{sed}. Essas ferramentas formam a base do
fluxo de trabalho de um desenvolvedor em ambiente Unix, permitindo versionar código,
automatizar tarefas repetitivas e transformar texto de forma programática, como
descrito em \cite{gitscm2024commandline}.

\section{Conceitos}

\begin{itemize}
  \item \textbf{Git}: sistema de controle de versão distribuído que registra o
    histórico de alterações de um projeto em commits, permitindo trabalho paralelo
    através de branches e a posterior integração via merge.
  \item \textbf{Make}: ferramenta de automação de compilação que, a partir de um
    \texttt{Makefile}, define regras (targets, dependências e comandos) para evitar
    a repetição manual de comandos como \texttt{pdflatex} ou \texttt{g++}
    \cite{embarcados2024makefile}.
  \item \textbf{Sed}: editor de fluxo (\emph{stream editor}) que aplica transformações
    de texto linha a linha, muito usado para substituições rápidas via expressões
    regulares sem abrir um editor interativo.
\end{itemize}

\bibliographystyle{plain}
\bibliography{referencias}

\end{document}
